\section*{Problem 4}
The Lucas numbers satisfy $L_n=L_{n-1}+L_{n-2}$ with $L_0=2$ and $L_1=1$.

\subsection*{(a) Show that $L_n=F_{n-1}+F_{n+1}$ for $n\ge2$, where $F_n$ is the $n$-th Fibonacci number.}
\textbf{Proof by induction on $n$.}
For $n=2$: $L_2=3$ and $F_{1}+F_{3}=1+2=3$.
Assume for some $k\ge2$ that $L_k=F_{k-1}+F_{k+1}$ and $L_{k-1}=F_{k-2}+F_k$.
Then
\[
L_{k+1}=L_k+L_{k-1}=(F_{k-1}+F_{k+1})+(F_{k-2}+F_k)
=(F_{k-2}+F_{k-1})+(F_k+F_{k+1})=F_k+F_{k+2},
\]
which is the desired form with $n=k+1$. Hence $L_n=F_{n-1}+F_{n+1}$ for all $n\ge2$. \qed

\subsection*{(b) Find an explicit formula for the Lucas numbers.}
The characteristic equation of $L_n=L_{n-1}+L_{n-2}$ is $x^2=x+1$ with roots
$\varphi=\frac{1+\sqrt5}{2}$ and $\psi=\frac{1-\sqrt5}{2}$. Solving
$L_n=A\varphi^n+B\psi^n$ using $L_0=2$ and $L_1=1$ gives $A=B=1$. Thus
\[
\boxed{\,L_n=\varphi^n+\psi^n\,}.
\]
